% ==============================================================================
% preamble/packages.tex — แพ็กเกจมาตรฐานสำหรับตำราวิชาการขนาด 6 x 9 นิ้ว
% ตามมาตรฐาน thai-academic-latex-formatting
% ==============================================================================

% 1. กำหนดขนาดรูปเล่ม 6 x 9 นิ้ว (Royal Octavo / Standard Academic Book)
\usepackage[
    paperwidth=6in,
    paperheight=9in,
    top=0.85in,
    bottom=0.85in,
    inner=1.00in,     % เผื่อเย็บเล่มด้านใน (Binding Gutter)
    outer=0.75in,
    footskip=0.45in,
    headheight=14pt
]{geometry}

% 2. ระบบฟอนต์และภาษาสำหรับ XeLaTeX
\usepackage{fontspec}
\usepackage{xunicode}
\usepackage{xltxtra}

% 3. การควบคุมระยะบรรทัดและการเยื้องย่อหน้า
\usepackage{setspace}
\usepackage{indentfirst} % บังคับย่อหน้าแรกหลังหัวข้อ (Heading) ให้เยื้องเท่ากับย่อหน้าปกติ

% 4. การจัดการหัวข้อและป้องกันหัวข้อตกค้าง (Heading Pagination)
\usepackage{titlesec}
\usepackage{needspace}

% 5. คำบรรยายภาพและตาราง (Caption Formatting)
\usepackage{caption}

% 6. รายการและโครงสร้าง (Lists)
\usepackage{enumitem}

% 7. ตารางและคณิตศาสตร์
\usepackage{booktabs}
\usepackage{tabularx}
\usepackage{amsmath}
\usepackage{amssymb}
\usepackage{array}

% 8. กราฟิกและกล่องข้อความ
\usepackage{graphicx}
\usepackage{xcolor}
\usepackage{tikz}
\usepackage[most]{tcolorbox}

% สภาพแวดล้อมกล่องข้อความวิชาการ (Academic Callout Box)
\newtcolorbox{academicNote}[1][]{
    enhanced,
    breakable,
    colback=teal!6!white,
    colframe=teal!75!black,
    fonttitle=\bfseries,
    title=ข้อสังเกตและแนวทางปฏิบัติ,
    arc=3mm,
    boxrule=0.8pt,
    left=12pt, right=12pt, top=10pt, bottom=10pt,
    #1
}

% 9. หัวกระดาษและท้ายกระดาษ (Running Headers)
\usepackage{fancyhdr}

% 10. ลิงก์และบุ๊กมาร์ก (Hyperref & xurl)
\usepackage{xurl}
\usepackage[
    colorlinks=true,
    linkcolor=black,
    citecolor=black,
    urlcolor=blue!70!black,
    bookmarksnumbered=true,
    pdfstartview=FitH
]{hyperref}
